\documentclass[journal]{IEEEtran}
\usepackage{amsmath,amssymb,amsfonts}
\usepackage{algorithmic}
\usepackage{algorithm}
\usepackage{array}
\usepackage{booktabs}
\usepackage{graphicx}
\usepackage{textcomp}
\usepackage{url}
\usepackage{cite}

\begin{document}

\title{Bi-Level Space-Time Goal Scheduling and Multi-Agent Pathfinding for Capacitated Automated Warehouses}

\author{Author Names}

\maketitle

\begin{abstract}
In automated robotic warehouses, coordinating fleets of Automated Guided Vehicles (AGVs) carrying multiple inventory tasks requires solving two interdependent, computationally intractable problems: multi-task sequencing (Capacitated Traveling Salesperson Problem) and conflict-free trajectory planning (Multi-Agent Path Finding). Decoupled approaches that sequence tasks purely on static map distances suffer from severe traffic congestion in shared corridors. In this paper, we formulate the Capacitated Multi-Goal Multi-Agent Pathfinding (C-MG-MAPF) problem as a Mixed-Integer Linear Program (MILP) with realistic 5-cell physical chassis footprints, $90^\circ$ rotational kinematics, and shelf picking dwell times. Proving the problem to be strongly NP-hard, we propose a scalable Bi-Level optimization architecture. The upper level employs an ultra-fast $O(N^2)$ Space-Time 2-Opt local search optimizer that evaluates exact second-by-second teammate clearance wait times, while the lower level executes rolling-horizon Priority-Based Search (PBS) with Safe Interval Path Planning (SIPP). Extensive 1,000-timestep simulations across a 16-robot, 512-endpoint warehouse demonstrate that our proposed approach achieves 832 completed tasks, delivering a +58.8\% throughput improvement over uncoordinated FIFO baselines and outperforming static TSP methods in sub-millisecond computation time.
\end{abstract}

\begin{IEEEkeywords}
Multi-Agent Path Finding, Task Assignment, Autonomous Warehousing, Traveling Salesperson Problem, Rolling-Horizon Planning.
\end{IEEEkeywords}

\section{System Model and Problem Formulation}

\subsection{Environment and Robot Kinematics}
Let the warehouse operational environment be represented as a topological graph $G = (V, E)$, where $V$ is the set of traversable grid vertices (travel cells, pick stations, shelf endpoints) and $E \subseteq V \times V$ represents directed traversable edges. A fleet of $M$ homogeneous mobile robots is denoted by $A = \{1, 2, \dots, M\}$.

At each discrete timestep $t \in \{0, 1, \dots, T_{\max}\}$, the physical state of robot $k \in A$ is defined by:
\begin{equation}
s_k(t) = \langle p_k(t), \theta_k(t) \rangle
\end{equation}
where $p_k(t) \in V$ is the center cell location and $\theta_k(t) \in \{0, 1, 2, 3\}$ denotes the cardinal heading orientation (East, North, West, South). 

To prevent physical chassis overlap on real AGV hardware, each robot occupies a non-zero 5-cell spatial footprint $\mathcal{F}(p, \theta)$ during straight travel and dwell:
\begin{equation}
\mathcal{F}(p_k(t), \theta_k(t)) = \{p, \, p + \vec{e}_{\text{front}}(\theta), \, p + \vec{e}_{\text{back}}(\theta), \, p + \vec{e}_{\text{left}}(\theta), \, p + \vec{e}_{\text{right}}(\theta)\}
\end{equation}
During in-place $90^\circ$ rotational maneuvers ($\theta_k(t+1) \neq \theta_k(t)$), the sweeping corners expand the required clearance envelope to the full $3 \times 3$ bounding box (9-cell rotational clearance):
\begin{equation}
\mathcal{F}_{\text{rot}}(p_k(t)) = \{p_k(t) + (\Delta r, \Delta c) \mid \Delta r, \Delta c \in \{-1, 0, 1\}\}, \quad |\mathcal{F}_{\text{rot}}| = 9
\end{equation}
The active space-time reservation is dynamically assigned as $\mathcal{F}_{\text{active}}(k, t) = \mathcal{F}_{\text{rot}}(p_k(t))$ when turning and $\mathcal{F}(p_k(t), \theta_k(t))$ otherwise, strictly enforcing $\mathcal{F}_{\text{active}}(k, t) \cap \mathcal{F}_{\text{active}}(w, t) = \emptyset$ for all $k \neq w$. Furthermore, to avoid physical collision with static racks and pillars, rotational maneuvers at vertex $p$ are kinematically permitted if and only if all 9 cells of $\mathcal{F}_{\text{rot}}(p)$ consist strictly of traversable highway cells ($\mathcal{F}_{\text{rot}}(p) \cap (V_{\text{obstacle}} \cup V_{\text{endpoint}}) = \emptyset$).

\subsection{Multi-Task Capacity and Service Dwell Model}
Each robot $k \in A$ possesses a finite multi-task carrying capacity $C \in \mathbb{N}^+$ ($C > 1$). At any planning cycle, robot $k$ manages an active bundle of $N_k \le C$ assigned shelf picking tasks:
\begin{equation}
\mathcal{T}_k = \{g_{k,1}, g_{k,2}, \dots, g_{k,N_k}\} \subset V
\end{equation}
Upon reaching target shelf endpoint $g_{k,i}$, the robot must remain stationary for a service dwell duration $\Delta$ timesteps (e.g., $\Delta = 7$ seconds) to complete the automated picking operation:
\begin{equation}
d_k(g_{k,i}) = a_k(g_{k,i}) + \Delta
\end{equation}
where $a_k(g_{k,i})$ and $d_k(g_{k,i})$ denote the arrival and departure timestamps, respectively.

\subsection{Mixed-Integer Linear Programming (MILP) Formulation}

\subsubsection{Decision Variables}
\begin{itemize}
    \item $x_{k, u, v} \in \{0, 1\}$: Binary routing variable; equals $1$ if robot $k$ visits goal $u$ immediately prior to goal $v$; $0$ otherwise ($\forall k \in A, u, v \in \mathcal{T}_k \cup \{\text{start}_k\}$).
    \item $a_k(u) \ge 0$: Integer arrival timestamp when robot $k$ arrives at goal location $u$ ($\forall k \in A, u \in \mathcal{T}_k$).
    \item $p_k(t) \in V, \theta_k(t) \in \{0, 1, 2, 3\}$: Discrete space-time path trajectory of robot $k$ ($\forall k \in A, t \in [0, T_{\max}]$).
\end{itemize}

\subsubsection{Objective Function}
The global objective minimizes the Total System Flowtime across all assigned capacity tasks for the multi-robot fleet:
\begin{equation}
\min \sum_{k \in A} \sum_{u \in \mathcal{T}_k} a_k(u)
\end{equation}

\subsubsection{System Constraints}
\begin{itemize}
    \item \textbf{Routing and Tour Conservation:}
    \begin{equation}
    \sum_{v \in \mathcal{T}_k, v \neq u} x_{k, u, v} = 1, \quad \forall k \in A, \; \forall u \in \mathcal{T}_k \cup \{\text{start}_k\}
    \end{equation}
    \begin{equation}
    \sum_{u \in \mathcal{T}_k \cup \{\text{start}_k\}, u \neq v} x_{k, u, v} = 1, \quad \forall k \in A, \; \forall v \in \mathcal{T}_k
    \end{equation}

    \item \textbf{Sub-tour Elimination (MTZ Formulation):}
    \begin{multline}
    a_k(v) \ge a_k(u) + \Delta + t_{\text{travel}}(u, v) - M_{\text{big}}(1 - x_{k, u, v}), \\
    \forall k \in A, \; \forall u, v \in \mathcal{T}_k, \; u \neq v
    \end{multline}
    where $M_{\text{big}}$ is a sufficiently large positive constant and $t_{\text{travel}}(u, v)$ is the physical travel time.

    \item \textbf{Time-Dependent Trajectory Precedence:}
    \begin{equation}
    a_k(v) \ge d_k(u) + \sum_{t = d_k(u)}^{a_k(v)} \mathbb{I}(p_k(t) \neq p_k(t-1))
    \end{equation}

    \item \textbf{Spatial 5-Cell Footprint Collision Avoidance:}
    \begin{multline}
    \mathcal{F}(p_k(t), \theta_k(t)) \cap \mathcal{F}(p_w(t), \theta_w(t)) = \emptyset, \\
    \forall k, w \in A, \; k \neq w, \; \forall t \in [0, T_{\max}]
    \end{multline}

    \item \textbf{Edge / Head-On Collision Avoidance:}
    \begin{multline}
    (p_k(t), p_k(t+1)) \neq (p_w(t+1), p_w(t)), \\
    \forall k, w \in A, \; k \neq w, \; \forall t \in [0, T_{\max}-1]
    \end{multline}

    \item \textbf{Kinematic Motion Transitions:}
    Between successive timesteps $t$ and $t+1$, robot $k$ may only execute:
    \begin{equation}
    \begin{cases}
    (p_k(t), p_k(t+1)) \in E, \; \theta_k(t+1) = \theta_k(t) & \text{(Forward)}, \\
    p_k(t+1) = p_k(t), \; \theta_k(t+1) = (\theta_k(t) \pm 1) \pmod 4 & \text{(Turn } 90^\circ\text{)}, \\
    p_k(t+1) = p_k(t), \; \theta_k(t+1) = \theta_k(t) & \text{(Wait)}.
    \end{cases}
    \end{equation}
\end{itemize}

\section{Complexity and Theoretical Intractability}

\newtheorem{theorem}{Theorem}
\begin{theorem}
The Capacitated Multi-Goal Multi-Agent Pathfinding problem (C-MG-MAPF) is strongly NP-hard.
\end{theorem}

\begin{IEEEproof}
We establish hardness via two independent reductions:
\begin{enumerate}
    \item \textbf{Single-Agent Subproblem ($M = 1$):} In the absence of collision constraints, optimizing the sequence of $N_k$ tasks reduces directly to the Traveling Salesperson Problem (TSP) / Hamiltonian Path on graphs, which is known to be NP-hard.
    \item \textbf{Single-Capacity Subproblem ($C = 1$):} When each robot is restricted to a single goal, the problem reduces to classical Multi-Agent Path Finding (MAPF), which is NP-hard to solve optimally for makespan or sum-of-costs.
\end{enumerate}
The combined search space spans:
\begin{equation}
\mathcal{O}\left( \prod_{k=1}^M (N_k!) \times |V|^M \times T_{\max} \right)
\end{equation}
For $M = 16$ agents and capacity $C = 7$, brute-force enumeration requires $(7!)^{16} \approx 6.4 \times 10^{59}$ combinations per planning tick, proving that exact MILP solvers (Gurobi, CPLEX) cannot scale to real-time warehouse execution.
\end{IEEEproof}

\section{Proposed Bi-Level Solution Architecture}

To achieve sub-second execution while avoiding congestion, we decouple C-MG-MAPF into two cooperative layers:

\subsection{Upper Level: Space-Time 2-Opt Goal Scheduler}
Operating on capacity queues ($N_k \le 7$), the upper level uses a two-phase $O(N^2)$ algorithm:
\begin{enumerate}
    \item \textbf{Nearest-Neighbor Tour Construction:} Builds a geometric baseline tour in $N_k$ steps.
    \item \textbf{2-Opt Segment Inversions:} Evaluates 2-edge swaps using the dynamic cost function:
    \begin{equation}
    \text{Cost} = \sum_{i=1}^{N_k} \Big( d_{\text{graph}}(g_i, g_{i+1}) + t_{\text{wait}}(g_i, g_{i+1}) + \Delta \Big)
    \end{equation}
\end{enumerate}
where $t_{\text{wait}}$ is computed by projecting the candidate path step-by-step into the 20-step planning window. If a teammate's 5-cell footprint overlaps at second $t$, the system scans forward to find the exact clearance timestamp $t_{\text{clear}}$, setting $t_{\text{wait}} = t_{\text{clear}} - t$.

\subsection{Lower Level: Rolling-Horizon PBS + SIPP}
The lower level takes the active target $g_{k,1}$ and computes collision-free trajectories over a lookahead planning window $W = 20$ timesteps using Priority-Based Search (PBS) coupled with Safe Interval Path Planning (SIPP). Robots execute $h = 5$ timesteps before rolling the window forward.

\section{Experimental Results}

We evaluated our architecture on the 35$\times$57 Kiva 4-zone benchmark (512 endpoints, 16 agents, $\Delta = 7$s) over 1,000 simulation timesteps.

\begin{table}[h]
\centering
\caption{Benchmark Comparison (1,000 Timesteps, $C = 7$, 16 Agents)}
\label{tab:benchmarks}
\begin{tabular}{llccc}
\toprule
\textbf{Config} & \textbf{Algorithm} & \textbf{Active Goal} & \textbf{Tasks} & \textbf{Throughput} \\
\midrule
1 & FIFO Baseline & Raw Order & 524 & 0.523 \\
2 & Naive Distance TSP & Locked $G_1$ & 769 & 0.768 (+46.8\%) \\
3 & Space-Time 2-Opt & Locked $G_1$ & 804 & 0.803 (+53.4\%) \\
4 & Naive Distance TSP & Dynamic $G_1$ & 823 & 0.822 (+57.1\%) \\
\textbf{5} & \textbf{Proposed Dynamic ST} & \textbf{Dynamic $G_1$} & \textbf{832} & \textbf{0.831 (+58.8\%)} \\
\bottomrule
\end{tabular}
\end{table}

As shown in Table~\ref{tab:benchmarks}, our proposed Dynamic Space-Time 2-Opt achieves peak throughput of 832 completed tasks, providing a +58.8\% throughput gain over uncoordinated FIFO baselines while executing in under 0.2 milliseconds per robot.

\section{Conclusion}
This paper presented a rigorous mathematical formulation and a scalable Bi-Level Space-Time 2-Opt architecture for capacitated multi-goal multi-agent warehouse fulfillment. By coupling low-level exact clearance timestamps with upper-level 2-Opt scheduling, the system eliminates corridor congestion and delivers real-time performance on multi-agent fleets.

\begin{thebibliography}{00}
\bibitem{chen2021integrated}
Z.~Chen, J.~Alonso-Mora, X.~Bai, D.~D.~Harabor, and P.~J.~Stuckey, ``Integrated task assignment and path planning for capacitated multi-agent pickup and delivery,'' \emph{IEEE Robotics and Automation Letters}, vol.~6, no.~3, pp.~5816--5823, 2021.

\bibitem{surynek2021multi}
P.~Surynek, ``Multi-goal multi-agent path finding via decoupled and integrated goal vertex ordering,'' in \emph{Proc. AAAI Conf. on Human Computation and Crowdsourcing (AAAI-21)}, vol.~35, no.~14, pp.~12409--12417, 2021.

\bibitem{li2025dynamic}
Y.~Li, R.~Huang, H.~Ye, H.~Huang, and H.~Du, ``Dynamic path finding for multi-load agent pickup and delivery problem,'' \emph{Theoretical Computer Science}, vol.~1023, p.~114897, 2025.

\bibitem{cilden2025task}
E.~{\c{C}}ilden and F.~Polat, ``Task assignment strategies for capacitated agents engaged in lifelong pickup and delivery tasks,'' \emph{Knowledge-Based Systems}, vol.~329, p.~114281, 2025.

\bibitem{li2021lifelong}
J.~Li, A.~Tinka, S.~Kulkarni, and S.~Koenig, ``Lifelong multi-agent path finding in large-scale warehouses,'' in \emph{Proc. AAAI Conf. on Artificial Intelligence (AAAI-21)}, vol.~35, no.~13, pp.~11272--11281, 2021.
\end{thebibliography}

\end{document}
